%% ======================================================================
%% 10_conclusions.tex -- 9. Conclusions
%% ======================================================================
\section{Conclusions}
\label{sec:conclusions}
AeroTray demonstrates that three independent balancing loops can be stacked on a mobile base using low-cost hardware, and that base-disturbance propagation across them can be both measured with synchronized logs and predicted from a first-principles model. Symbolically verified and parameterized from the CAD assembly, datasheets, and logs, the model yields six experimentally confirmed findings.

(i)~The aerial module behaves as an inverted pendulum on an accelerating pivot whose center of mass, identified from hover telemetry, lies $0.14\pm\SI{0.06}{\metre}$ above the spherical joint and \SI{9.5}{\milli\metre} toward M8. (ii)~The attitude loop works only because proportional stiffness exceeds gravitational stiffness by a narrow margin supplied by the underestimated thrust gain; that narrow margin explains the slow attitude mode observed in hover. (iii)~The slow divergences in early flights were a torque-authority limit, with a region of attraction of only 3--4$^\circ$, and were resolved by widening the angular-rate command limit rather than retuning gains. (iv)~The ball loop is bounded by its delay and by a right-half-plane zero caused by ball height above the plate pivot, which together set the derivative-gain ceiling; effective plant gain is roughly 70\% of the CAD value due to mechanism backlash. (v)~The passive spherical joint removes slope from the ball dynamics, turns are benign, and residual ball error with spinning rotors is dominated by a slowly varying direct aerodynamic and vibrational load rather than octocopter attitude motion.

(vi)~The drone is justified on the ramp for two quantified reasons. By leveling itself, it decouples the slope from the plate and preserves 84--95\% of servo travel, whereas a base-fixed tray would exhaust its travel on an \SI{11}{\degree} slope. By supporting part of its own weight, it acts as adjustable ballast that lowers the effective center of mass: with the rotors carrying half the module weight, the uphill tip-over angle rises from \SI{21.5}{\degree} (measured on the \SI{10.4}{\kilo\gram} full robot) to a predicted \SI{27.7}{\degree}, the downhill tip-over angle rises from \SI{30.5}{\degree} to a predicted \SI{38.2}{\degree}, and the acceleration margin at \SI{20}{\degree} increases by a factor of 3.5. On the ramp, the drone should not hover; instead, it should operate in a 35--55\% weight-relief window at a cost of 2--5\% battery capacity per minute. Conversely, hover thrust margin depends on LiPo voltage squared: at 12.4\,V the module lifts off at 1725--\SI{1775}{\micro\second}, whereas at 11.1\,V it would require 1850--\SI{2000}{\micro\second}, explaining why hover failed at low charge and limiting useful flight to 3--\SI{7}{\minute} per battery pair.

Experimentally, the final firmware holds the \SI{3.6}{\kilo\gram} module on the vertical guide within 0.55--\SI{1.7}{\degree} RMS attitude error, and the ball controller keeps the ball within a median of 19--\SI{30}{\milli\metre} of center on a stationary base, on flat ground, and on a \SI{10}{\degree} ramp. In the final full-robot trials (Nov.~23), 13 balancing trials lasted 23--\SI{122}{\second} (\SI{816}{\second} total; 29--\SI{44}{\milli\metre} median), 10 of them with spinning rotors and the drone leveled at 0.71--\SI{1.0}{\degree} RMS. In four full-system trials, the base crossed the \SI{10}{\degree} ramp with rotors spinning and the ball balancing for 18--\SI{30}{\second} (30--\SI{35}{\milli\metre} median, drone at 0.85--\SI{1.15}{\degree} RMS with zero cutoffs, the ball exiting the plate in the final segment). The module also took off and landed on the vertical guide carrying the tray, though the ball rolled off at liftoff. Base speed measured that day agrees with the model within 8--10\%. Two concrete fixes follow directly from these trials: re-leveling the plate's $x$-axis neutral, which concentrates the error, and smoothing guide takeoff, whose stiction jump dislodged the ball at liftoff. The model also identifies the highest-impact next steps: measuring battery voltage to compensate collective, testing ramp mode, lowering the module center of mass below the spherical joint, shortening ball-loop delay, shielding the plate from rotor airflow, and smoothing base velocity references. All design files, source code, analysis scripts, videos, and raw data are published so these experiments can be reproduced and extended---for example, using the already implemented tilt feedforward or controllers trained on the published model.
